\documentclass[12pt, a4paper]{article}

% --- PACOTES BÁSICOS E IDIOMA ---
\usepackage[utf8]{inputenc}      % Suporte a caracteres UTF-8
\usepackage[brazil]{babel}       % Traduz termos e ajusta hifenização
\usepackage[T1]{fontenc}         % Suporte a acentuação correta
\usepackage{indentfirst}         % Indenta o primeiro parágrafo
\usepackage{xcolor}              % Suporte a cores
\usepackage{fancyhdr}            % Personalização de cabeçalhos e rodapés
\usepackage{graphicx}            % Inclusão de imagens/logos
\usepackage{helvet}              % Pacote para fonte estilo Arial/Helvetica
\usepackage{enumitem}            % Personalização avançada de listas/itens
\usepackage{amsmath}             % Suporte a símbolos e comandos matemáticos (\text)
\usepackage{calc} % Habilita o uso de \widthof para cálculos de dimensão
\usepackage{booktabs}
\usepackage{siunitx}
\usepackage{circuitikz}
\usepackage{float}
\usepackage{pgfplots}
\usepgfplotslibrary{groupplots}
\usepackage{titlesec}

\titleformat{\section}
  {\normalfont\large\bfseries}
  {\thesubsection}
  {1em}
  {}

\titleformat{\subsection}
  {\normalfont\normalsize\bfseries}
  {\thesubsection}
  {1em}
  {}

\pgfplotsset{compat=1.18}

% --- SELEÇÃO DE FONTE DO DOCUMENTO ---
% Descomente a linha da fonte que deseja utilizar no documento inteiro:
\usepackage{mathptmx}            % Opção 1: Times New Roman
%\usepackage{helvet} \renewcommand{\familydefault}{\sfdefault} % Opção 2: Arial

% --- FORMATAÇÃO DE PÁGINA E MARGENS ---
\usepackage{geometry}
\geometry{
    a4paper,
    top=4.2cm,                   % Ajustado para acomodar o cabeçalho de 1.5cm
    bottom=2.0cm,
    left=3.0cm,
    right=2.0cm,
    headheight=1.5cm,            % Altura do cabeçalho igual à da logo (1.5cm)
    headsep=0.8cm
}

% --- DEFINIÇÃO DE CORES ---
\definecolor{laranjaufc}{HTML}{E66100} % Tom de laranja institucional
\definecolor{cinzaresposta}{HTML}{2B2B2B} % Tom escuro e elegante para a resposta/debate

% --- CONFIGURAÇÃO DO CABEÇALHO ---
\pagestyle{fancy}
\fancyhf{} % Limpa cabeçalho e rodapé padrão

\fancyhead[C]{%
    \setlength{\fboxsep}{0pt}% Remove espaçamento interno padrão do colorbox
    % 1. Armazena a imagem em uma caixa temporária para medir sua largura exata
    \sbox0{\includegraphics[height=1.5cm]{UFC_logo.png}}%
    \vbox to 1.5cm{%
        \vfil
        \hbox to \textwidth{%
            % 2. Exibe a logo alinhada à esquerda
            \usebox0%
            % 3. Espaço de 1 cm entre a logo e o retângulo
            \hspace{0.5cm}%
            % 4. Retângulo laranja preenchendo o restante da linha até a margem direita
            \colorbox{laranjaufc}{%
                \makebox[\dimexpr\textwidth-\wd0-1cm\relax][l]{%
                    \vbox to 1.5cm{%
                        \vfil
                        \hbox{%
                            \hspace{0.25cm}% Margem de 0.2cm da extremidade esquerda do retângulo
                            \fontfamily{phv}\selectfont % Define fonte Arial (Helvetica)
                            \color{white}\bfseries\small
                            \begin{tabular}{@{}l@{}}
                                UNIVERSIDADE FEDERAL DO CEARÁ -- UFC \\
                                DEPARTAMENTO DE ENGENHARIA ELÉTRICA
                            \end{tabular}%
                        }%
                        \vfil
                    }%
                }%
            }%
        }%
        \vfil
    }%
}

% Remove a linha horizontal padrão do fancyhdr
\renewcommand{\headrulewidth}{0pt}

% Número da página no rodapé à direita
\fancyfoot[R]{\thepage}

% Remove a linha horizontal padrão do fancyhdr
\renewcommand{\headrulewidth}{0pt}

% --- VARIÁVEIS DE IDENTIFICAÇÃO DO RELATÓRIO ---
\newcommand{\disciplina}{Laboratório de Eletrônica de Potência}
\newcommand{\coddisciplina}{TH258}
\newcommand{\professor}{Prof. Ernande Eugênio}
\newcommand{\numpratica}{Documento Digital - Projeto de um conversor não-isolado CC-CC Buck}